\documentclass{article}
\usepackage{/globals/de}
\usepackage{/globals/anki}
\ankiprefix{wahlen}
\ankitopic{Wahlen}
\ankishow
\begin{document}
\section{Wahlen}
Wer eine Wahl gewonnen hat, und ob es überhaupt eine GewinnerIn gibt, kann auf mehrere Art und Weisen definiert sein; je nach dem nach welcher \emph{Art} der Mehrheit gesucht wird.
\begin{description}
 \item[Relative Mehrheit] Die Option mit den meisten Stimmen gewinnt.
 \item[Einfache Mehrheit] Die Option, wenn es sie gibt, die mehr als $50\%$ der \emph{abgegebenen} Stimmen hat gewinnt.
 \item[Absolute Mehrheit] Die Option, wenn es sie gibt, die mehr als $50\%$ der \emph{abgebbaren} Stimmen hat gewinnt.
 \item[Qualifizierte Mehrheit] Die letzten beiden, aber die Grenze selbst liegt bei über $50\%$. So kann nach \zb $2/3$ der abgegebenen oder abgebbaren Stimmen gefragt sein.
\end{description}
Ungültige Stimmen und Enthaltungen zählen nicht als Stimmen.
\ankinote{mehrheit-rel}{Mehrheit Relativ}{Meiste Stimmen}
\ankinote{mehrheit-einf}{Mehrheit Einfach}{Über 50% der Stimmen}
\ankinote{mehrheit-abs}{Mehrheit Absolut}{Über 50% der Personen, nicht Stimmen}
\ankinote{mehrheit-qual}{Mehrheit Qualifiziert}{Einfach oder Absolut, aber mit mehr als 50p}

\subsection{Beschlussfähigkeit}
Wird der Bundestag als Beispiel genommen, so kann angemerkt werden dass oftmals auch bei relativen und einfachen Mehrheiten eine bestimmte Mehrheit der abgebbaren Stimmen benötigt wird, weil der Bundestag nur \emph{beschlussfähig} ist, also \zb Gesetze verabschieden kann, wenn über $50\%$ der Mitglieder anwesend sind.
\ankinote{beschlussfähigkeit}{BT Beschlussfähigkeit}{Über 50% anwesend}
\end{document}